\documentclass[10pt, aspectratio=169]{beamer}
\usepackage{../beamer_theme_408}

\title{408考研零基础 --- 计算机组成原理}
\subtitle{2-10 虚拟存储器}
\author{}
\date{}

\begin{document}


\begin{frame}{本节目标}
    \begin{enumerate}
        \item 理解虚拟存储器的基本概念与作用
        \item 掌握分页存储管理（页、页框、页表、地址转换）
        \item 理解分段存储管理的基本思想
        \item 了解段页式存储管理
    \end{enumerate}
\end{frame}

\begin{frame}[t]{为什么需要虚拟存储器？}
    \begin{block}{面临的问题}
        \begin{enumerate}
            \item \textbf{物理内存有限}：运行程序可能大于物理内存
            \item \textbf{多道程序}：多个程序同时运行，总需求 > 物理内存
            \item \textbf{地址空间隔离}：不同程序不应相互干扰
        \end{enumerate}
    \end{block}
    \vspace{0.3em}
    \begin{block}{虚拟存储器的解决方案}
        \begin{itemize}
            \item 给每个程序提供\textbf{独立的虚拟地址空间}
            \item 运行时再将虚拟地址转换为物理地址
            \item 部分内容可在磁盘上，需要时调入
        \end{itemize}
    \end{block}
    \vspace{0.3em}
    \begin{center}
        \begin{tabular}{|c|c|}
            \hline
            \textbf{虚拟地址（VA）} & 程序看到的地址 \\
            \hline
            \textbf{物理地址（PA）} & 主存中的真实地址 \\
            \hline
        \end{tabular}
    \end{center}
\end{frame}

\begin{frame}{虚拟存储器与Cache的类比}
    \begin{center}
        \begin{tabular}{|c|c|}
            \hline
            \textbf{Cache} & \textbf{虚拟存储器} \\
            \hline
            主存 $\to$ Cache & 磁盘 $\to$ 主存 \\
            \hline
            块（Block/Line） & 页（Page） \\
            \hline
            标记（Tag） & 页号 \\
            \hline
            命中（Hit） & 页命中（Page in Memory） \\
            \hline
            缺失（Miss） & 缺页（Page Fault） \\
            \hline
        \end{tabular}
    \end{center}
    \vspace{0.5em}
    \begin{itemize}
        \item 虚拟存储器是\textbf{主存-磁盘}层次的Cache
        \item 页大小通常为4KB（比Cache行大得多）
        \item 缺页处理时间极长（毫秒级）——尽量提高命中率
    \end{itemize}
\end{frame}

\begin{frame}[t]{分页存储管理——基本概念}
    \begin{block}{核心概念}
        \begin{itemize}
            \item \textbf{页（Page）}：虚拟地址空间划分的固定大小块
            \item \textbf{页框（Page Frame）}：物理内存划分的固定大小块
            \item \textbf{页大小相同}——通常4KB（$2^{12}$）
        \end{itemize}
    \end{block}
    \vspace{0.3em}
    \begin{block}{虚拟地址结构}
        \[
            \underbrace{\text{虚拟页号（VPN）}}_{\text{高位}} \quad
            \underbrace{\text{页内偏移（Offset）}}_{\text{低位}}
        \]
    \end{block}
    \vspace{0.3em}
    \begin{exampleblock}{例：虚拟地址32位，页大小4KB}
        \begin{align*}
            \text{页内偏移} &= \log_2(4K) = 12\text{位} \\
            \text{虚拟页号} &= 32 - 12 = 20\text{位} \\
            \text{虚拟地址空间} &= 2^{20} \text{页} = 1M\text{页}
        \end{align*}
    \end{exampleblock}
\end{frame}

\begin{frame}{页表}
    \begin{block}{页表的作用}
        \begin{itemize}
            \item 存储\textbf{虚拟页号到物理页框号的映射关系}
            \item 每个进程有\textbf{独立的页表}
        \end{itemize}
    \end{block}
    \vspace{0.3em}
    \begin{center}
        \begin{tabular}{|c|c|c|}
            \hline
            \textbf{虚拟页号} & \textbf{物理页框号} & \textbf{状态位} \\
            \hline
            0 & 5 & 有效 \\
            \hline
            1 & -- & 无效（在磁盘上） \\
            \hline
            2 & 9 & 有效 \\
            \hline
            $\cdots$ & $\cdots$ & $\cdots$ \\
            \hline
        \end{tabular}
    \end{center}
    \vspace{0.3em}
    \begin{itemize}
        \item \textbf{状态位}（存在位）：该页是否在主存中
        \item 若无效 $\to$ 缺页中断 $\to$ 从磁盘调入
    \end{itemize}
\end{frame}

\begin{frame}{地址转换过程——分页}
    \begin{center}
        \begin{tabular}{c}
            CPU给出虚拟地址（VPN + Offset） \\
            $\downarrow$ \\
            以VPN为索引访问\textbf{页表} \\
            $\downarrow$ \\
            \begin{tabular}{c|c}
                页表项有效（页在主存） & 页表项无效（缺页） \\
                $\downarrow$ & $\downarrow$ \\
                读PPN（物理页框号） & 缺页中断 \\
                $\downarrow$ & $\downarrow$ \\
                物理地址 = PPN + Offset & 从磁盘读入页框 \\
                $\downarrow$ & $\downarrow$ \\
                访问主存 & 更新页表后重试 \\
            \end{tabular}
        \end{tabular}
    \end{center}
    \vspace{0.3em}
    \begin{center}
        \textcolor{blue}{每次地址转换都需要查页表——效率是关键}
    \end{center}
\end{frame}

\begin{frame}{TLB——快表}
    \begin{block}{TLB (Translation Lookaside Buffer)}
        \begin{itemize}
            \item 为解决页表查询速度慢而引入
            \item 是页表的\textbf{硬件Cache}（位于CPU内部）
            \item 存储最近使用的页表项
        \end{itemize}
    \end{block}
    \vspace{0.3em}
    \begin{center}
        \begin{tabular}{c}
            CPU给出VPN \\
            $\downarrow$ \\
            \begin{tabular}{c|c}
                TLB命中 & TLB缺失 \\
                $\downarrow$ & $\downarrow$ \\
                直接得到PPN & 查页表 \\
                （1个时钟周期） & $\downarrow$ \\
                              & 更新TLB \\
            \end{tabular}
        \end{tabular}
    \end{center}
    \vspace{0.3em}
    \begin{itemize}
        \item TLB通常为全相联或高度组相联
    \end{itemize}
\end{frame}

\begin{frame}[t]{分段存储管理}
    \begin{block}{基本概念}
        \begin{itemize}
            \item 按\textbf{逻辑单位}（程序段、数据段、栈段等）划分
            \item 每段有\textbf{段名/段号}和\textbf{段长}
            \item 段长可以不固定
        \end{itemize}
    \end{block}
    \vspace{0.3em}
    \begin{block}{逻辑地址结构}
        \[
            \underbrace{\text{段号}}_{\text{找到段基址}} \quad
            \underbrace{\text{段内偏移}}_{\text{需小于段长}}
        \]
    \end{block}
    \vspace{0.3em}
    \begin{block}{段表}
        \begin{center}
            \begin{tabular}{|c|c|c|}
                \hline
                \textbf{段号} & \textbf{段长} & \textbf{段基址} \\
                \hline
                0 & 64KB & 0x100000 \\
                \hline
                1 & 32KB & 0x200000 \\
                \hline
            \end{tabular}
        \end{center}
    \end{block}
\end{frame}

\begin{frame}[t]{分段 vs 分页}
    \begin{center}
        \begin{tabular}{|c|c|c|}
            \hline
            \textbf{特性} & \textbf{分页} & \textbf{分段} \\
            \hline
            划分依据 & 固定大小（硬件） & 逻辑单位（用户/编译器） \\
            \hline
            大小是否固定 & 固定（如4KB） & 不固定 \\
            \hline
            地址空间 & 线性一维 & 二维（段号+偏移） \\
            \hline
            共享/保护 & 较困难 & 容易（按段） \\
            \hline
            碎片问题 & 内部碎片 & 外部碎片 \\
            \hline
        \end{tabular}
    \end{center}
    \vspace{0.5em}
    \begin{itemize}
        \item 分页：关注\textbf{内存利用率}，对程序员透明
        \item 分段：关注\textbf{逻辑组织}，对程序员可见
    \end{itemize}
\end{frame}

\begin{frame}{段页式存储管理}
    \begin{block}{基本思想——分段+分页的结合}
        \begin{enumerate}
            \item 先将程序按逻辑分段（段式管理）
            \item 每段内再分成固定大小的页（页式管理）
            \item 物理内存按页框管理
        \end{enumerate}
    \end{block}
    \vspace{0.3em}
    \begin{block}{逻辑地址结构}
        \[
            \underbrace{\text{段号}}_{\text{查段表}} \quad
            \underbrace{\text{段内页号}}_{\text{查页表}} \quad
            \underbrace{\text{页内偏移}}_{\text{块内地址}}
        \]
    \end{block}
    \vspace{0.3em}
    \begin{block}{地址转换过程}
        \begin{itemize}
            \item 段号 $\to$ 段表 $\to$ 得到该段页表基址
            \item 段内页号 $\to$ 页表 $\to$ 得到物理页框号
            \item 页框号 + 页内偏移 $\to$ 物理地址
        \end{itemize}
    \end{block}
\end{frame}

\begin{frame}{段页式地址转换流程}
    \begin{center}
        \begin{tabular}{c}
            逻辑地址（段号 $|$ 段内页号 $|$ 偏移） \\
            $\downarrow$ \\
            查\textbf{段表}：段号 $\to$ 该段页表起始地址 \\
            $\downarrow$ \\
            查\textbf{页表}：段内页号 $\to$ 物理页框号 \\
            $\downarrow$ \\
            物理地址 = 页框号 $\times$ 页大小 + 偏移 \\
            $\downarrow$ \\
            访问主存
        \end{tabular}
    \end{center}
    \vspace{0.5em}
    \begin{itemize}
        \item \textcolor{red}{两次访存查表}（段表+页表）$\to$ 速度慢
        \item 改进：使用TLB缓存最近使用的段表/页表项
    \end{itemize}
\end{frame}

\begin{frame}{本节总结}
    \begin{enumerate}
        \item 虚拟存储器：给每个程序提供独立虚拟地址空间
        \item 分页：固定大小，通过页表+TLB地址转换
        \item 分段：按逻辑单位，段表记录基址和段长
        \item 段页式：分段+分页结合，两次查表
    \end{enumerate}
\end{frame}

\end{document}
